% =====================================================================
\chapter{Первое включение и настройки пульта}
\label{ch:radio}
% =====================================================================
\chapterintro
  {какие предупреждения появляются при включении, как откалибровать стики, выбрать режим и
   порядок каналов, настроить звук, голос, подсветку, контроль батареи и оборудование}
  {пульт с актуальной SD-картой}

\section{Предупреждения при включении}

EdgeTX заботится о безопасности и при включении проверяет положение органов управления:
\begin{description}[font=\sffamily\bfseries\color{etx}]
  \item[Throttle warning] стик газа не на минимуме. Опустите его --- предупреждение исчезнет.
  \item[Switch warning] переключатели не в «стартовых» положениях, заданных для модели.
        На экране показано, какие переключатели нужно перевести.
  \item[SD card / Storage warning] версия файлов на SD-карте не совпадает с прошивкой
        (глава~\ref{ch:firmware}).
  \item[Alarms disabled / Sound off] звук выключен --- легко пропустить важный сигнал.
\end{description}

Любое предупреждение можно пропустить кнопкой, но делать так по привычке не стоит: однажды
модель с включённым армом и газом не на нуле рванёт с места у вас в руках.

\section{Калибровка стиков и крутилок}

Калибровка --- первое, что нужно сделать на новом пульте, после замены стиков или прошивки.

\begin{enumerate}
  \item \mpath{SYS\sep Hardware\sep Calibration} (в некоторых версиях --- отдельная страница
        \menu{Calibration}).
  \item Нажмите \btn{ENTER} --- \textbf{отпустите стики в центр}, крутилки поставьте в среднее
        положение и снова \btn{ENTER}.
  \item Несколько раз \textbf{обведите стики по кругу до упора} и прокрутите крутилки
        от края до края.
  \item Нажмите \btn{ENTER} для сохранения.
\end{enumerate}

\begin{tip}
Проверьте результат на странице с положениями стиков главного экрана или в
\mpath{MDL\sep Outputs}: в центре должно быть 0, в крайних положениях --- ровно $\pm100$.
Если газ в минимуме показывает, например, $-97$, калибровка сделана плохо --- повторите.
\end{tip}

\section{Radio Setup (общие настройки)}

\mpath{SYS\sep Radio Setup} --- длинная страница. Разберём пункты по группам.

\subsection{Дата, время, единицы}

\menu{Date}, \menu{Time} --- нужны для правильных имён файлов логов. \menu{Units} ---
метрические (\menu{Metric}) или имперские. \menu{Time zone} и \menu{Adjust RTC} позволяют
подстраивать часы по GPS из телеметрии.

\subsection{Батарея пульта}

\begin{center}\small
\begin{tabularx}{\linewidth}{@{}p{36mm}X@{}}
\toprule
\menu{Battery range} & шкала индикатора батареи на главном экране. Для 2S (две 18650 или 2S LiPo)
разумно \code{6.6}--\code{8.4}\,В. \\
\menu{Battery low} & порог предупреждения. Для 2S Li-ion/LiPo --- \code{6.6}--\code{7.0}\,В. \\
\menu{Battery calib.} & подстройка показаний: измерьте напряжение мультиметром и выставьте
такое же на экране. \\
\bottomrule
\end{tabularx}
\end{center}

\subsection{Звук, голос, вибрация}

\begin{itemize}
  \item \menu{Sound mode} (режим звука): \menu{Beeper}, \menu{All} (всё включено) и
        промежуточные. Для голосовых сообщений нужен режим, в котором разрешён голос.
  \item \menu{Volume}, \menu{Beep volume}, \menu{Wav volume}, \menu{Background volume} ---
        громкость общего звука, писков, голосовых файлов и фоновой музыки.
  \item \menu{Voice language} --- язык голосовых сообщений. Должна существовать папка
        \code{SOUNDS/<язык>} на SD-карте (например, \code{SOUNDS/ru} или \code{SOUNDS/en}).
  \item \menu{Haptic} --- вибромотор: режим, сила, длительность.
  \item \menu{Vario} --- тон «пищалки» вариометра для планеров.
\end{itemize}

\subsection{Предупреждения (Alarms)}

\menu{Inactivity alarm} напомнит о забытом включённом пульте (например, через 10 минут
без движения стиков). \menu{Sounds off} предупреждает о выключенном звуке.
\menu{Memory low}/\menu{Storage} --- о заполненной карте.

\subsection{Подсветка и экран}

\menu{Backlight mode} (\menu{Keys}, \menu{Controls}, \menu{Keys \& Ctrls}, \menu{On}),
\menu{Duration} --- время работы подсветки, \menu{Brightness}, \menu{Contrast}.
Подсветка --- заметный потребитель энергии; разумно выставить 10--30 секунд.

\subsection{Режим стиков и порядок каналов}

\begin{description}[font=\sffamily\bfseries\color{etx}]
  \item[Mode] режим стиков 1--4 (глава~1). Обычно \textbf{2}.
  \item[Default channel order] порядок, в котором стики попадают в каналы \sw{CH1}--\sw{CH4}
        \emph{при создании новой модели}. Буквы: \textbf{A}ileron (крен), \textbf{E}levator
        (тангаж), \textbf{T}hrottle (газ), \textbf{R}udder (рыскание).
        \begin{itemize}
          \item \textbf{AETR} --- стандарт для ExpressLRS/Crossfire и Betaflight/INAV
                (\sw{CH1} крен, \sw{CH2} тангаж, \sw{CH3} газ, \sw{CH4} рыскание);
          \item \textbf{TAER} --- порядок Spektrum;
          \item для FrSky-приёмников в самолётах часто используют AETR.
        \end{itemize}
\end{description}

\begin{note}
\menu{Default channel order} влияет \emph{только на новые модели}. У существующих моделей
порядок меняется в миксерах. Полётному контроллеру (Betaflight, INAV, ArduPilot) важно, чтобы
его «карта каналов» (channel map) совпадала с тем, что реально передаёт пульт.
\end{note}

\subsection{USB}

\menu{USB mode}: \menu{Ask} (спрашивать при подключении), \menu{Joystick} (пульт как
игровой контроллер --- для симуляторов), \menu{Storage} (флешка), \menu{Serial} (виртуальный
COM-порт для отладки и прошивки модуля ELRS). Удобнее всего оставить \menu{Ask}.

\subsection{Прочее}

\menu{Splash screen} --- время показа заставки. \menu{Play delay} --- пауза между
голосовыми сообщениями о переключателях. \menu{Model quick select} --- быстрый выбор модели
долгим нажатием. \menu{Owner} (имя владельца) --- полезно для учебного класса,
где много одинаковых пультов.

\section{Hardware (оборудование)}

\mpath{SYS\sep Hardware} описывает физическую конфигурацию пульта:

\begin{itemize}
  \item \textbf{Stick/pot names and types}: имена стиков и крутилок (можно переименовать
        \sw{S1} в \sw{Cam}), тип крутилки --- с фиксацией в центре или без, либо \menu{None}.
  \item \textbf{Switches}: тип каждого переключателя --- \menu{None}, \menu{Toggle}
        (без фиксации, моментальный), \menu{2POS}, \menu{3POS}; имя переключателя.
  \item \textbf{ADC filter} --- сглаживание аналоговых сигналов (оставьте включённым,
        кроме особых случаев).
  \item \textbf{Internal RF}: тип внутреннего модуля (\menu{CRSF} для ELRS, \menu{Multi}
        для 4-в-1) и скорость связи с ним.
  \item \textbf{Serial port}: назначение последовательных портов (AUX): выдача телеметрии
        наружу, подключение внешнего GPS, отладка, «беспроводной тренер».
  \item \textbf{Bluetooth}, \textbf{Sticks/gyro} --- если есть в вашей сборке.
\end{itemize}

\begin{caution}
Если в \menu{Hardware} тип переключателя указан неверно (например, двухпозиционный вместо
трёхпозиционного), среднее положение не будет распознаваться. Если тип внутреннего модуля
выбран неправильно, пульт не сможет общаться с модулем, и привязка не пройдёт.
\end{caution}

\section{Global Functions и Trainer}

\mpath{SYS\sep Global Functions} --- специальные функции, общие для \emph{всех} моделей
(например, объявлять голосом напряжение пульта или делать скриншоты). Модель использует их,
если в её настройках разрешено \menu{Use global functions}. Подробно специальные функции ---
в главе~\ref{ch:special}.

\mpath{SYS\sep Trainer} --- настройка режима «учитель--ученик», глава~\ref{ch:trainer}.

\begin{tasks}
\begin{enumerate}
  \item Выполните калибровку и убедитесь, что все стики дают 0 в центре и $\pm100$ по краям.
  \item Установите \menu{Battery range} и \menu{Battery low} для своего типа аккумулятора.
  \item Проверьте типы всех переключателей в \menu{Hardware}, пощёлкав ими.
  \item Установите \menu{Default channel order} = AETR и \menu{Voice language} = ru
        (если установлен русский голосовой пакет).
\end{enumerate}
\end{tasks}
